\begin{abstract}
Multilingual Retrieval-Augmented Generation (mRAG) systems frequently retrieve documents in multiple languages simultaneously. Because web corpora vary in coverage and recency across languages, high-resource English passages and lower-resource non-English passages can assert mutually exclusive factual claims. While prior studies examine language bias in retrieval, reranking, or citation selection over complementary passages, how large language models (LLMs) resolve direct cross-lingual factual contradictions remains unexplored. We introduce a counterbalanced evaluation protocol on a 150-item English--Korean benchmark derived from MKQA (\textsc{MKQA-Conflict-150}), spanning 50 named entities, 50 dates, and 50 numeric facts ($2{,}100$ inference runs per model). For each query, we pair ground-truth claim $A$ against a type-matched counterfactual $B$ across swapped language assignments ($\text{Pair}_1$: $[\text{EN}{=}A, \text{KO}{=}B]$ vs.\ $\text{Pair}_2$: $[\text{KO}{=}A, \text{EN}{=}B]$), crossed with query language and passage order. Evaluating \texttt{Qwen-2.5-7B-Instruct} and \texttt{Llama-3.1-8B-Instruct} via exact extraction and concept-marginalized log-probabilities reveals a strong English authority bias: overall English Preference Rate (EPR) reaches \qwenEPRAll\% for Qwen and \llamaEPRAll\% for Llama ($p \qwenMcNemarPAll$). Crucially, when the English passage asserts a falsehood against a truthful Korean passage ($\text{Pair}_2$), models override factual truth in \qwenPairTwoOverrideAll\% and \llamaPairTwoOverrideAll\% of cases, even when the fact is known parametrically.
\end{abstract}
